\section{Обоснование и описание архитектуры и стека технологий}

\subsection{Обоснование выбора архитектуры}

\subsubsection{Анализ альтернативных архитектур}

При проектировании системы интеллектуального поиска и генерации ответов на основе документов рассматривалось несколько архитектурных подходов, каждый из которых имеет свои сильные и слабые стороны.

Первый из рассмотренных подходов — монолитная архитектура, которая предполагает размещение всей логики приложения в единой кодовой базе и единственном процессе. На первый взгляд, такой подход привлекателен простотой начальной разработки и отладки. Однако применительно к решаемой задаче монолит обладает несколькими существенными недостатками. Прежде всего, он не позволяет независимо масштабировать ресурсоёмкие компоненты — в частности, сервис векторизации текста, требующий значительных вычислительных мощностей, оказывается привязан к остальной части приложения. Кроме того, жёсткая зависимость компонентов друг от друга серьёзно затрудняет замену или обновление отдельных частей системы: например, смена модели эмбеддингов или языковой модели может потребовать пересборки и переразвёртывания всего приложения целиком. Наконец, в монолитной архитектуре отказ одного компонента с высокой вероятностью приводит к недоступности всей системы, что критично для промышленного использования.

Альтернативой выступает микросервисная архитектура, которая устраняет перечисленные недостатки за счёт декомпозиции системы на независимые сервисы. Каждый такой сервис выполняет строго определённую функцию, а взаимодействие между ними осуществляется через чётко очерченные интерфейсы, как правило, поверх REST API. Такой подход обеспечивает слабую связанность компонентов и открывает возможности для гибкого управления системой.

С учётом детального анализа требований к разрабатываемой системе в качестве основного архитектурного решения была выбрана именно микросервисная архитектура. Этот выбор обусловлен несколькими ключевыми критериями. Во-первых, это масштабируемость: микросервисный подход допускает горизонтальное масштабирование сервиса векторизации и сервиса поиска по мере роста нагрузки, без необходимости расширять всю систему целиком. Во-вторых, модульность: микросервисная архитектура позволяет заменить языковую модель или модель эмбеддингов практически без изменений в остальных компонентах, что особенно важно в быстро меняющейся области NLP. В-третьих, отказоустойчивость: сбои изолируются в рамках одного сервиса и не приводят к коллапсу всей системы. Наконец, гибкость развёртывания: отдельные компоненты можно разворачивать на специализированном оборудовании — например, GPU-ускоренный сервис эмбеддингов может работать на выделенных вычислительных узлах, тогда как лёгковесные сервисы пользовательского интерфейса могут обходиться без ускорителей.

\subsubsection{Обоснование выбора паттерна RAG}

Для решения задачи интеллектуального поиска и ответов на вопросы по корпоративным документам в качестве основного архитектурного паттерна был выбран RAG (Retrieval-Augmented Generation). Этот выбор обусловлен ограничениями альтернативных подходов и спецификой решаемой задачи.

Чисто генеративный подход, то есть использование языковой модели без дополнительного контекста, оказался неприменим в данном случае по нескольким причинам. Языковые модели обладают фиксированными знаниями, ограниченными датой их обучения, и поэтому попросту не осведомлены о внутренних документах организации, которые составляют основу предметной области. Более того, генерация ответов без опоры на реальные документы неизбежно приводит к галлюцинациям — явлению, при котором модель выдаёт правдоподобные на вид, но фактически неверные или выдуманные сведения. Альтернативный путь — дообучение (fine-tuning) модели на корпоративных данных — также имеет серьёзные недостатки: он требует значительных вычислительных ресурсов, занимает много времени и, что самое важное, не позволяет оперативно обновлять базу знаний при появлении новых документов или изменении старых.

Паттерн RAG решает все перечисленные проблемы принципиально иным способом. Вместо того чтобы полагаться на зафиксированные в весах модели знания, RAG-подход даёт модели релевантные фрагменты документов непосредственно в контексте запроса. Это обеспечивает фактическую точность ответов, поскольку модель опирается на реальные, только что извлечённые данные. Кроме того, обновление базы знаний становится простой и быстрой операцией — достаточно добавить новые документы в индекс, и модель сразу начнёт их учитывать, без какого-либо переобучения. Именно сочетание этих свойств — точности, актуальности и лёгкости обновления — сделало RAG архитектурной основой разрабатываемой системы.


\subsubsection{Обоснование разделения на конвейеры Ingest и Retrieval}

В основе архитектуры разрабатываемой системы лежит разделение процессов обработки данных на два независимых конвейера: конвейер индексации (Ingest Pipeline) и конвейер поиска (Retrieval Pipeline). Такое разделение продиктовано принципиально разными требованиями к этим процессам и необходимостью обеспечить плавную работу системы под реальной нагрузкой.

Конвейер индексации запускается асинхронно в тот момент, когда в систему загружаются новые документы. Этот процесс включает в себя парсинг исходных файлов, разбиение текста на смысловые фрагменты (чанки) и векторизацию — то есть преобразование каждого фрагмента в многомерный вектор-эмбеддинг. Все эти операции требуют значительных вычислительных ресурсов и могут занимать заметное время, особенно при работе с большими объёмами документов. Именно поэтому конвейер индексации вынесен в отдельный сервис с поддержкой очереди задач. Пользователи при этом никак не блокируются: они могут продолжать отправлять запросы, даже когда в фоне загружается и индексируется новый пакет документов.

Конвейер поиска, напротив, работает в синхронном режиме и запускается непосредственно в ответ на каждый запрос пользователя. Его задача — максимально быстро найти релевантные фрагменты в уже проиндексированной базе знаний и сформировать ответ. Время отклика здесь является критическим показателем, поэтому конвейер поиска должен быть максимально лёгким и оптимизированным.

Разделение двух конвейеров даёт ключевое преимущество: каждый из них можно масштабировать и настраивать независимо, в соответствии с характером его нагрузки. Например, при резком увеличении потока пользовательских запросов можно горизонтально масштабировать именно сервисы поиска, не трогая ресурсоёмкие компоненты индексации. И наоборот, если в организацию поступает большой объём новых документов, можно выделить дополнительные вычислительные мощности для конвейера индексации, не опасаясь замедлить ответы текущим пользователям. Такая архитектурная гибкость особенно важна для корпоративных систем с неравномерной и трудно предсказуемой нагрузкой.


\subsection{Обоснование выбора технологий и инструментов}

\subsubsection{Выбор языка программирования — Python}

В качестве основного языка программирования для всех микросервисов разрабатываемой системы был выбран Python. Этот выбор обусловлен несколькими весомыми причинами, главная из которых — положение Python как де-факто стандарта в области машинного обучения и обработки естественного языка. Благодаря этому разработчики получают доступ к богатейшей экосистеме готовых библиотек, покрывающих практически любую мыслимую задачу. Особенно важна широкая поддержка со стороны ключевых фреймворков для работы с большими языковыми моделями — таких как LangChain, LlamaIndex и HuggingFace Transformers. Кроме того, Python известен высокой скоростью разработки прототипов и возможностью итеративно улучшать систему без необходимости переписывать большие фрагменты кода при смене подходов. Для проекта, в котором многие архитектурные решения принимаются экспериментально, это свойство оказывается критически важным.

\subsubsection{Выбор фреймворка FastAPI}

Для реализации REST API каждого микросервиса в качестве основного фреймворка выбран FastAPI. Этот выбор продиктован сочетанием производительности, удобства разработки и богатой встроенной функциональности.

FastAPI обладает высокой производительностью, поскольку он построен на ASGI-сервере Starlette и поддерживает асинхронную обработку запросов. Для решаемой задачи это свойство особенно ценно: большинство операций в системе являются I/O-интенсивными — обращения к базе данных, вызовы внешних API, работа с векторным хранилищем — и асинхронный подход позволяет эффективно утилизировать ресурсы, не блокируя потоки в ожидании ответа.

Ещё одним важным преимуществом FastAPI является автоматическая генерация документации в форматах Swagger UI и ReDoc. Эта документация создаётся прямо на основе аннотаций типов Python, что практически исключает риск расхождения между реализацией API и его описанием. Для системы, состоящей из нескольких микросервисов, которые активно взаимодействуют друг с другом, такая автоматическая документация заметно упрощает интеграцию и снижает время на согласование интерфейсов.

Кроме того, FastAPI предоставляет встроенную валидацию данных через библиотеку Pydantic. Это позволяет на уровне фреймворка проверять корректность входящих и исходящих данных, что снижает вероятность ошибок при передаче информации между сервисами и избавляет разработчика от необходимости писать валидацию вручную.

Для полноты обоснования стоит сравнить FastAPI с основным альтернативным вариантом — фреймворком Flask. Flask, будучи более старым и проверенным решением, работает в синхронном режиме WSGI и не обеспечивает нативной поддержки асинхронности. В ситуациях с высокой нагрузкой и большим количеством I/O-операций Flask значительно уступает FastAPI в производительности. Более того, Flask не имеет встроенной системы валидации данных, что вынуждает разработчика подключать сторонние библиотеки вроде Marshmallow или писать проверки вручную. Совокупность этих факторов — асинхронность, производительность, автоматическая документация и встроенная валидация — сделала FastAPI наиболее подходящим выбором для разрабатываемой системы.


\subsubsection{Векторная база данных — ChromaDB}

Для хранения векторных представлений текстовых фрагментов (эмбеддингов) в качестве основной базы данных выбрана ChromaDB. Этот выбор обусловлен, прежде всего, простотой интеграции: ChromaDB легко разворачивается как отдельный сервис и предоставляет как Python SDK, так и полноценный REST API, что хорошо сочетается с выбранной микросервисной архитектурой. Важной особенностью ChromaDB является поддержка хранения метаданных непосредственно вместе с векторами — это позволяет хранить исходный текст фрагмента (чанка) и всю метаинформацию о документе (название, автор, дату загрузки и так далее) в едином хранилище, без необходимости в отдельной базе данных для метаданных. Кроме того, ChromaDB предоставляет встроенную поддержку различных метрик сходства, включая косинусное и евклидово расстояние, что даёт гибкость при настройке алгоритмов поиска. Немаловажным фактором является также открытый исходный код и активная поддержка со стороны сообщества.

При выборе векторной базы данных рассматривались и альтернативные решения, в частности Qdrant и Weaviate. Оба этих инструмента обладают богатой функциональностью и высокой производительностью, однако требуют заметно более сложной настройки и сопровождения. ChromaDB была выбрана как наиболее простое в развёртывании и интеграции решение, которое при этом полностью покрывает потребности проекта на этапе разработки и способно справляться с небольшими производственными нагрузками. Сочетание достаточной функциональности и минимального порога входа сделало ChromaDB оптимальным выбором.

\subsubsection{Реляционная база данных — PostgreSQL}

Для хранения структурированных данных — метаданных документов, истории диалогов, информации о пользователях и настроек системы — в качестве основной реляционной базы данных выбрана СУБД PostgreSQL. Этот выбор объясняется в первую очередь надёжностью и зрелостью технологии: PostgreSQL имеет многолетнюю историю использования в самых разных производственных системах, от небольших проектов до крупных корпоративных инсталляций, и зарекомендовала себя как исключительно стабильное решение.

Отдельного внимания заслуживает богатая поддержка типов данных в PostgreSQL. В частности, наличие типа JSON позволяет гибко хранить метаданные произвольной структуры без необходимости заранее проектировать жёсткую схему таблиц. Это свойство оказывается очень полезным, поскольку разные типы документов могут требовать разного набора метаполей. Помимо этого, PostgreSQL обладает развитой экосистемой инструментов для резервного копирования и мониторинга, что упрощает эксплуатацию системы. И наконец, поддержка ACID-транзакций гарантирует целостность и согласованность данных даже в сценариях с параллельным доступом, что критически важно для сохранения истории диалогов и пользовательских данных.

\subsubsection{Объектное хранилище — MinIO}

Для хранения исходных файлов документов, загружаемых пользователями в систему, применяется объектное хранилище MinIO. Ключевым преимуществом MinIO является его полная совместимость с Amazon S3 API. Это означает, что при необходимости можно заменить MinIO на облачное хранилище S3 от любого провайдера без единого изменения в коде системы — достаточно поменять настройки подключения. Такая гибкость оставляет пространство для будущей миграции без архитектурной перестройки.

MinIO допускает самостоятельное развёртывание в Docker-контейнере, что полностью соответствует требованию автономности системы: все компоненты могут работать на собственной инфраструктуре предприятия, без передачи данных во внешние облака. Хранилище поддерживает файлы произвольного размера и формата, что важно, поскольку загружаемые документы могут быть самыми разными — от небольших текстовых файлов до объёмных PDF-отчётов. Кроме того, MinIO показывает высокую производительность при работе с бинарными данными.

Стоит подчеркнуть, что хранение файлов в объектном хранилище, а не в файловой системе сервера или в реляционной базе данных, является архитектурно правильным решением. Оно обеспечивает масштабируемость — объектное хранилище легко расширяется добавлением новых узлов — и независимость хранилища файлов от вычислительных узлов. При таком подходе можно перезапустить или заменить сервер, обрабатывающий запросы пользователя, без потери доступа к загруженным документам, что повышает отказоустойчивость системы в целом.


\subsubsection{Очередь задач — Redis Queue (RQ)}

Для асинхронной обработки задач индексации документов в системе применяется Redis в связке с библиотекой RQ (Redis Queue). Выбор именно этой комбинации обусловлен стремлением найти баланс между функциональностью и простотой. RQ заметно проще в настройке и использовании по сравнению с более тяжёлыми альтернативами, такими как Apache Kafka, RabbitMQ или даже Celery, при этом предоставляя всю необходимую функциональность для решаемой задачи. В данной конфигурации Redis используется исключительно как брокер задач, что позволяет не вводить в архитектуру лишней технологической сложности. Дополнительным преимуществом является возможность мониторинга состояния очереди через встроенный дашборд RQ.

Главная функциональная причина использования очереди задач заключается в том, что операции загрузки и обработки документов (парсинг, разбиение на чанки, генерация эмбеддингов и индексация) являются длительными и ресурсоёмкими. Если бы система выполняла их синхронно, в рамках HTTP-запроса пользователя, время ожидания ответа было бы совершенно неприемлемым. Очередь задач решает эту проблему кардинально: сервис немедленно возвращает пользователю ответ о том, что документ принят в обработку, а сам процесс индексации выполняется в фоновом режиме. Такой подход сохраняет отзывчивость системы даже при загрузке больших объёмов документов.

\subsubsection{Модели эмбеддингов — HuggingFace Transformers}

Для вычисления векторных представлений текста (эмбеддингов) используется экосистема HuggingFace Transformers. Этот выбор даёт доступ к широчайшему каталогу предобученных моделей различного размера и специализации, что позволяет экспериментировать и подбирать оптимальный вариант под конкретную задачу и доступные вычислительные ресурсы. Особенно важно, что в каталоге присутствуют мультиязычные модели, например \texttt{intfloat/multilingual-e5}, которые хорошо работают с русскоязычными документами — а это прямое требование предметной области.

Ключевым преимуществом использования локальных моделей эмбеддингов из HuggingFace является возможность их выполнения на собственной инфраструктуре, без передачи данных во внешние API. Это полностью соответствует общему требованию приватности и конфиденциальности корпоративных данных: тексты документов авиапредприятия никогда не покидают внутренний контур. Технически система организует кэширование весов модели в локальной директории \texttt{.hf\_cache}, что исключает повторную загрузку модели при каждом перезапуске контейнера и ускоряет развёртывание.

\subsubsection{Контейнеризация — Docker и Docker Compose}

Вся разрабатываемая система развёртывается с использованием Docker и Docker Compose. Выбор этой технологии контейнеризации продиктован несколькими взаимосвязанными причинами. Прежде всего, это воспроизводимость среды: образы контейнеров фиксируют все зависимости — от версии операционной системы до конкретных версий библиотек Python — что полностью исключает проблему «работает на моей машине» при переносе системы между средой разработки, тестирования и продуктивным окружением.

Кроме того, Docker обеспечивает изоляцию сервисов: каждый микросервис выполняется в своём изолированном окружении, что исключает конфликты зависимостей между компонентами. Например, разные сервисы могут использовать разные версии одной и той же библиотеки, и это не приведёт к ошибкам.